% Section 5.6, from lectures/L05/appendix/b_exercises.md.

\section{Exercises}
\label{c5:sec:exercises}

\begin{aside}
\textbf{How to work these.} Each exercise is tagged with the kinds of work it asks for:
\kinds{Paper} by hand, \kinds{CircuitVerse} in the browser, \kinds{VHDL} in a module of your own.

\textbf{How to check your work.} Every exercise that asks you to write a VHDL module ships a
self-checking testbench under \file{lectures/L05/exercises} in the companion repository. Write your
module in its \file{exercises/<module>} directory, using the entity name and the \textbf{port
order} the exercise specifies, then run it with GHDL; \secref{c2:app:c} has the three commands.

No FPGA board is needed for any exercise. Your job is to get the VHDL right, and the testbench is
how you confirm it. Short answers to the parts you work on paper are in
Appendix~\ref{app:answers}; reference solutions to the VHDL are published in the companion
repository, under \file{lectures/L05/solutions}. Try each exercise yourself before looking at
either.
\end{aside}

% ------------------------------------------------------------------------------------------
\exerciseset{\code{variable} Versus \code{signal}}
\label{c5:set:1}

\exercise{Tracing the signal accumulator}{Paper}
\label{c5:ex:1}
The broken \code{parity_gen} architecture from \secref{c5:a2} is reproduced below:

\begin{vhdlcode}
architecture broken of parity_gen is
signal parity_s: std_logic;
begin
    parity <= parity_s;

    process(bits) is
    begin
        parity_s <= '0';
        for i in bits'range loop
            parity_s <= parity_s xor bits(i);
        end loop;
    end process;
end architecture;
\end{vhdlcode}

\noindent Assume:
\begin{itemize}
  \item \code{bits} has the range \code{7 downto 0}.
  \item \code{bits = "10110110"}.
  \item \code{parity_s} holds \code{'0'} before the process begins.
\end{itemize}

\begin{parts}
  \item Trace the process by hand. For the initial assignment and every loop iteration, record:
    \begin{itemize}
      \item The current index \code{i}.
      \item The value of \code{bits(i)}.
      \item The value read from \code{parity_s}.
      \item The value scheduled for \code{parity_s}.
      \item Whether that scheduled assignment survives until the process suspends.
    \end{itemize}
    Remember:
    \begin{itemize}
      \item A signal assignment does not update the signal immediately.
      \item Assignments to the same signal during one process activation replace earlier scheduled
        assignments for the same simulation time.
    \end{itemize}
  \item Determine the final value assigned to \code{parity_s} after the process suspends. Then:
    \begin{itemize}
      \item Calculate the correct XOR parity of \code{"10110110"} by hand.
      \item Compare it with the result from the broken architecture.
      \item Explain why repeatedly assigning to the signal does not create an accumulator.
    \end{itemize}
  \item Trace the variable-based \code{behaviour} architecture from the chapter's
    \file{parity_gen.vhd} using the same input. For every loop iteration, record:
    \begin{itemize}
      \item The value of \code{bits(i)}.
      \item The value of \code{acc} before the XOR operation.
      \item The value of \code{acc} immediately after the XOR operation.
    \end{itemize}
    Confirm that:
    \begin{itemize}
      \item \code{acc} retains the result of each iteration for use by the next iteration.
      \item The final value assigned to \code{parity} is the correct XOR parity.
    \end{itemize}
\end{parts}

\exercise{The smaller of two values}{VHDL}
\label{c5:ex:2}
Write an entity named \code{min_of_two}.

\noindent The entity must have:

\begin{booktable}{@{}p{3em}p{4.9em}p{14.4em}L@{}}
  \thead{Port} & \thead{Direction} & \thead{Type} & \thead{Description} \\
  \midrule
  \code{a} & in & \code{std_logic_vector(3 downto 0)} & First value. \\
  \code{b} & in & \code{std_logic_vector(3 downto 0)} & Second value. \\
  \code{m} & out & \code{std_logic_vector(3 downto 0)} & The smaller of the two. \\
\end{booktable}

\noindent Treat \code{a} and \code{b} as unsigned four-bit values.

\noindent Implement the circuit using a process:
\begin{itemize}
  \item Include \code{a} and \code{b} in the sensitivity list.
  \item Declare a process-local variable to hold the smaller value.
  \item Compare \code{a} and \code{b} as unsigned numbers.
  \item Assign the variable exactly once during each process activation.
  \item Assign the variable to \code{m} after the comparison.
\end{itemize}
Use the facilities from \code{ieee.numeric_std} when performing the unsigned comparison, following
\secref{c2:a3}.

This exercise does not require a loop. Its purpose is to use a variable as temporary process-local
storage for something \emph{other than} an accumulator: \secref{c5:a2}'s worked example builds a
value up across the iterations of a loop, which is the case a variable is usually introduced for,
and it is easy to come away thinking that is the only thing they are good for. Here the variable
simply holds an intermediate result for the few lines between computing it and using it, which is
the more common use of the two in real code.

\modulefigure{min_of_two}

\selfcheck name your entity \code{min_of_two}, with ports \code{a}, \code{b} (in) and \code{m}
(out), all \code{std_logic_vector(3 downto 0)}, declared in that order; its testbench is in
\file{exercises/min_of_two}.

\exercise{Counting set bits}{VHDL, Paper}
\label{c5:ex:3}
Write an entity named \code{ones_count8} that counts how many bits of an 8-bit input are set.

Exercise~\ref{c5:ex:1} had you trace the accumulator bug on paper and Exercise~\ref{c5:ex:2} used a
variable that was not an accumulator at all. This one is the case \secref{c5:a2} is actually about:
a value built up across the iterations of a loop. It is the only exercise in the book where you
write one.

\noindent The entity must have:

\begin{booktable}{@{}p{3em}p{4.9em}p{14.4em}L@{}}
  \thead{Port} & \thead{Direction} & \thead{Type} & \thead{Description} \\
  \midrule
  \code{bits} & in & \code{std_logic_vector(7 downto 0)} & The value to count over. \\
  \code{ones} & out & \code{natural range 0 to 8} & The number of bits of \code{bits} that are
    \code{'1'}. \\
\end{booktable}

\noindent Implement it using a single combinational process:
\begin{itemize}
  \item Put \code{bits} in the sensitivity list.
  \item Declare a process-local \code{variable} for the running count, and set it to \code{0} at the
    top of every activation. A variable keeps its value between activations, so a count that is
    never cleared keeps growing.
  \item Loop over \code{bits'range}, adding \code{1} for each bit that is set.
  \item Assign the variable to \code{ones} once, after the loop.
\end{itemize}
Then answer, in a sentence each:
\begin{itemize}
  \item Write the same thing with a \code{signal} accumulator instead, run the testbench, and record
    what it reports for the input \code{00000010}. Explain the number you get using the scheduling
    rule, not by guessing.
  \item The loop runs eight times. How many adders does this design put on the FPGA, and how many
    clock cycles does it take to produce an answer? If those two answers surprise you, reread
    \secref{c5:loopnot}.
  \item \code{ones} is a \code{natural range 0 to 8} rather than a
    \code{std_logic_vector(3 downto 0)}. Both can hold the answer. What does the range buy you, and
    what would GHDL do if your count ever left it?
\end{itemize}

\modulefigure{ones_count8}

\selfcheck name your entity \code{ones_count8}, with input \code{bits}
(\code{std_logic_vector(7 downto 0)}) and output \code{ones} (\code{natural range 0 to 8}), declared
in that order; its testbench is in \file{exercises/ones_count8} and sweeps all 256 input values.

% ------------------------------------------------------------------------------------------
\exerciseset{A Module Worth Having on Its Own}
\label{c5:set:2}

\exercise{A synchronizer module}{VHDL, Paper}
\label{c5:ex:4}
Write the double-flop synchronizer as a module of its own, named \code{sync}.

You have already built this circuit. It is the first two flip-flops of Chapter~\ref{c4:ch}'s
\code{button_sync} (Exercise~\ref{c4:ex:9}), written inline there because at that point it was the
\emph{idea} that mattered: an asynchronous input needs two flip-flops before anything else may look
at it. Here you pull those two out into a module, and the exercise is not the logic, which you know,
but everything around it.

Two things make it worth the second pass. The first is that a bare synchronizer is what most
asynchronous inputs actually need: \code{button_sync} adds an edge detector because a button press
is an event, but a serial line, a status flag from another clock domain, or a switch you only ever
read the level of needs the two flip-flops and nothing more. Chapter~\ref{c8:ch}'s serial receiver
instantiates exactly this on its \code{rx} pin. The second is that it takes \textbf{two generics},
and one of them is not a number.

\begin{booktable}{@{}p{4.2em}p{10.8em}p{3.8em}L@{}}
  \thead{Generic} & \thead{Type} & \thead{Default} & \thead{Description} \\
  \midrule
  \code{SIZE} & \code{natural range 1 to 15} & \code{1} & How many independent signals to
    synchronize, and therefore the width of both vector ports. \\
  \code{PRESET} & \code{std_logic} & \code{'0'} & The value both stages take on reset. \\
\end{booktable}

\noindent and these ports:

\begin{booktable}{@{}p{5.6em}p{4.9em}>{\raggedright\arraybackslash}p{11.2em}L@{}}
  \thead{Port} & \thead{Direction} & \thead{Type} & \thead{Description} \\
  \midrule
  \code{clock} & in & \code{std_logic} & 50\,MHz system clock. \\
  \code{reset_s2_n} & in & \code{std_logic} & Active-low, \textbf{already synchronized} reset. A
    synchronizer does not synchronize its own reset; Chapter~\ref{c4:ch}'s \code{reset_sync} does
    that. \\
  \code{async_in} & in & \code{std_logic_vector(SIZE - 1 downto 0)} & The asynchronous input or
    inputs, straight from outside the clock domain. \\
  \code{sync_out} & out & \code{std_logic_vector(SIZE - 1 downto 0)} & The same signals, safe to
    use, two rising edges later. Each bit is synchronized independently. \\
\end{booktable}

\begin{parts}
  \item Implement it:
    \begin{itemize}
      \item Two internal \code{std_logic_vector(SIZE - 1 downto 0)} signals, one per stage.
      \item One process, sensitive to \code{clock} and \code{reset_s2_n}.
      \item On reset, drive both stages to \code{PRESET}, outside the clocked branch.
      \item On a rising clock edge, shift \code{async_in} through the two stages.
      \item Drive \code{sync_out} from the second stage with a concurrent assignment.
    \end{itemize}
  \item \code{PRESET} is a \code{std_logic} generic, which is the first one in this book that is not
    a number, and it exists because the safe reset value depends on what the signal \emph{means}.
    Only the instantiating design knows that. Answer:
    \begin{itemize}
      \item \code{button_sync} synchronizes active-low buttons, so it would pass \code{'1'}. What
        would \code{'0'} claim about the buttons for the two cycles after every reset?
      \item A serial line idles high, so Chapter~\ref{c8:ch}'s receiver passes \code{'1'} too. What
        would \code{'0'} look like to a receiver watching for a start bit?
      \item Given both, why is the \emph{default} \code{'0'} rather than \code{'1'}, and what does
        that say about who is responsible for getting it right?
    \end{itemize}
  \item This is also a question about what a generic costs. \code{SIZE} is
    \code{natural range 1 to 15} and \code{PRESET} is \code{std_logic}, and both are fixed before
    synthesis runs. Using \secref{c5:a3}'s picture of what the FPGA is actually made of, say how many
    flip-flops an instance at \code{SIZE = 4} uses, how many an instance at \code{SIZE = 1} uses,
    and where in the built circuit \code{PRESET} ends up. Neither generic appears as logic anywhere;
    say what happened to them instead.
  \item Verify it with its testbench. It instantiates your module twice, once one bit wide with
    \code{PRESET = '1'} and once three bits wide with \code{PRESET = '0'}, so a design that
    hard-codes either generic fails. It checks that \code{sync_out} follows \code{async_in} after
    \textbf{exactly} two rising edges rather than one or eventually, that each bit of a vector is
    synchronized on its own, and that reset loads \code{PRESET} with no clock edge involved.
\end{parts}
\note[Tip] write the architecture for \code{SIZE = 1} in your head first, then check that every line
is already correct for a vector. The aggregate \code{(others => PRESET)} and a plain vector
assignment both work element-wise, so the generalization should cost no extra code.

\modulefigure{sync}

\selfcheck name your entity \code{sync}, with generics \code{SIZE} (\code{natural range 1 to 15},
default \code{1}) and \code{PRESET} (\code{std_logic}, default \code{'0'}), inputs \code{clock},
\code{reset_s2_n}, \code{async_in} (\code{std_logic_vector(SIZE - 1 downto 0)}) and output
\code{sync_out} (\code{std_logic_vector(SIZE - 1 downto 0)}), declared in that order; its testbench
is in \file{exercises/sync}.

\textbf{Keep this file.} Chapter~\ref{c8:ch}'s serial receiver capstone instantiates it on the
\code{rx} pin, and asks you to copy it in.

% ------------------------------------------------------------------------------------------
\exerciseset{From VHDL to Hardware}
\label{c5:set:3}
The course demonstrates the Quartus workflow on a real DE0-CV board, first with the braking
assistant back in Chapter~\ref{c1:ch}, and Appendix~\ref{app:quartus} writes it up for reference.
You don't need Quartus or a board yourself; the exercises below are about understanding what that
workflow does to the VHDL you write.

\exercise{What Quartus does to your VHDL}{Paper}
\label{c5:ex:5}
Answer each of the following, referring to Appendix~\ref{app:quartus}. One or two sentences each.

\begin{parts}
  \item A Quartus project needs a \textbf{top-level entity} selected before it will compile:
    \begin{itemize}
      \item What does the tool do differently with the top-level entity than with any other entity
        in the project?
      \item \code{min_of_two} from Exercise~\ref{c5:ex:2} and its testbench are both valid VHDL.
        Why can only one of them ever be the top-level entity of an FPGA project?
    \end{itemize}
  \item Pin assignment connects a port name to a physical pin:
    \begin{itemize}
      \item Why must pins be assigned \emph{before} compiling, rather than afterwards?
      \item What would \code{m} physically be connected to if you compiled \code{min_of_two} having
        assigned no pins at all?
    \end{itemize}
  \item Synthesis turns your architecture into gates on the chip:
    \begin{itemize}
      \item Chapter~\ref{c2:ch}'s \code{combo_logic} exercise (Exercise~\ref{c2:ex:10}) had you
        write \code{x = ab + c'd} twice: once as one expression, and once with \code{c'd} named by
        an internal signal. What should synthesis produce for each, and why is your answer the same
        for both?
      \item Which of the two would you rather read six months from now?
    \end{itemize}
  \item Compilation produces warnings as well as errors:
    \begin{itemize}
      \item Why is ``it compiled with no errors'' a much weaker statement about an FPGA design than
        ``it compiled with no errors'' about a C program?
      \item Name one class of warning you would never ignore.
    \end{itemize}
\end{parts}

\exercise{Predicting parity}{Paper, VHDL}
\label{c5:ex:6}
\code{parity_gen} from \secref{c5:a2} is this chapter's worked example, and it ships a reference
testbench at \file{lectures/L05/parity_gen/parity_gen_tb.vhd}.

\begin{parts}
  \item Before running anything, predict \code{parity} by hand for each of these inputs:
    \begin{itemize}
      \item \code{"00000000"}
      \item \code{"10110110"}
      \item \code{"11111111"}
      \item \code{"10000000"}
    \end{itemize}
  \item Run the reference testbench over the worked example and confirm your four predictions:
\begin{shell}
cd lectures/L05/parity_gen
ghdl -a --std=93 parity_gen.vhd parity_gen_tb.vhd
ghdl -e --std=93 parity_gen_tb
ghdl -r --std=93 parity_gen_tb --assert-level=error --stop-time=10ms
\end{shell}
  \item Take any one of those inputs and flip exactly one bit, leaving the other seven unchanged:
    \begin{itemize}
      \item What happens to \code{parity}?
      \item Explain why changing exactly one input bit must \emph{always} invert the XOR-parity
        result, whichever bit you pick and whatever the starting value.
      \item This is the property that makes a parity bit able to detect a single-bit transmission
        error. What kind of error does it fail to detect, and why?
    \end{itemize}
\end{parts}

\exercise{Reviewing a latched decoder}{Paper}
\label{c5:ex:7}
A colleague sends you this combinational decoder for review. It compiles, and Quartus reports no
errors.

\begin{vhdlcode}
architecture behaviour of level_decode is
begin
    process(sel) is
    begin
        case sel is
            when "00"   => leds <= "0001";
            when "01"   => leds <= "0010";
            when "10"   => leds <= "0100";
            when others => null;
        end case;
    end process;
end architecture;
\end{vhdlcode}

\begin{parts}
  \item Quartus emits \code{Warning: Inferred latch(es) for signal "leds"}. Using \secref{c5:a3},
    explain what the tool built and why. What did the VHDL ask for that combinational logic cannot
    do?
  \item The author's defence is that \code{sel} is two bits, so \code{"11"} is the only value
    \code{when others} can match, and the design never generates it. Give two separate reasons
    that argument does not save them:
    \begin{itemize}
      \item one about what \code{sel} can actually hold, from \secref{c1:a4}.
      \item one that would still apply even if \code{sel} really could only ever be \code{"00"},
        \code{"01"} or \code{"10"}.
    \end{itemize}
  \item Fix it, in two different ways:
    \begin{itemize}
      \item by changing only the \code{when others} branch.
      \item by adding a single line \emph{before} the \code{case} statement, leaving every branch
        untouched.
    \end{itemize}
    Which of the two would you rather maintain, and why? The second scales to a \code{case} with
    twenty branches; the first does not.
  \item This is a warning, not an error, and the design will often appear to work on the board. Why
    does that make it more dangerous than an error rather than less?
\end{parts}
